\documentclass[11pt,a4paper]{article}

% --- Required Packages ---
\usepackage[margin=20mm,top=22mm,bottom=22mm]{geometry}
\usepackage{amsmath,amssymb}
\usepackage{graphicx}
\usepackage{xcolor}
\usepackage{fancyhdr}
\usepackage{listings}
\usepackage{tcolorbox}
\usepackage{array}
\usepackage{colortbl}
\usepackage{tikz}
\usetikzlibrary{calc}
\usepackage{eso-pic}
\usepackage[hidelinks]{hyperref}

% --- Color Definitions ---
\definecolor{navyblue}{RGB}{13,35,58}
\definecolor{codegreen}{rgb}{0,0.6,0}
\definecolor{codegray}{rgb}{0.5,0.5,0.5}
\definecolor{codepurple}{rgb}{0.58,0,0.82}
\definecolor{backcolour}{rgb}{0.96,0.97,0.98}

% --- Double Border Setup ---
\AddToShipoutPictureBG{%
  \begin{tikzpicture}[remember picture, overlay]
    \draw[line width=1.5pt, navyblue] 
      ($(current page.north west)+(1.0cm,-1.0cm)$) 
      rectangle 
      ($(current page.south east)+(-1.0cm,1.0cm)$);
    \draw[line width=0.5pt, navyblue] 
      ($(current page.north west)+(1.15cm,-1.15cm)$) 
      rectangle 
      ($(current page.south east)+(-1.15cm,1.15cm)$);
  \end{tikzpicture}%
}

% --- Code Listing Style ---
\lstdefinestyle{pythonstyle}{
    backgroundcolor=\color{backcolour},   
    commentstyle=\color{codegreen}\itshape,
    keywordstyle=\color{navyblue}\bfseries,
    numberstyle=\tiny\color{codegray},
    stringstyle=\color{codepurple},
    basicstyle=\ttfamily\small,
    breakatwhitespace=false,         
    breaklines=true,                 
    captionpos=b,                    
    keepspaces=true,                 
    numbers=left,                    
    numbersep=8pt,                  
    showspaces=false,                
    showstringspaces=false,
    showtabs=false,                  
    tabsize=4,
    frame=single,
    rulecolor=\color{navyblue!30}
}

\pagestyle{plain}

\begin{document}

% =========================================================
% FRONT PAGE (TITLE PAGE)
% =========================================================
\begin{titlepage}
\centering
\vspace*{0.5cm}

% Placeholder for Logo
\begin{tikzpicture}
  \node[draw=navyblue, line width=1.2pt, rounded corners=6pt, inner sep=14pt, fill=navyblue!5] {
    \begin{minipage}{0.7\textwidth}
      \centering
      \Huge \textbf{\textcolor{navyblue}{[ LOGO PLACEHOLDER ]}}\\[6pt]
      \small Department of Artificial Intelligence \& Machine Learning
    \end{minipage}
  };
\end{tikzpicture}

\vspace{1.2cm}

{\Huge \textbf{\textcolor{navyblue}{Albertian Institute of Science \& Technology}}}\\[0.4cm]
{\large (Approved by AICTE, Affiliated to State Board of Technical Education)}\\[1.8cm]

{\Huge \textbf{Lab Manual for}}\\[0.8cm]

{\Huge \textbf{\textcolor{navyblue}{Python Programming Lab -- 3347}}}\\[0.6cm]

{\Huge \textbf{Semester -- III (AI \& ML)}}\\[4.0cm]

\vfill

{\Large \textit{Prepared By,}}\\[0.3cm]
{\Large \textbf{Kalyani M Biju}}\\[0.2cm]
{\large Lecturer, Department of Artificial Intelligence \& Machine Learning}

\end{titlepage}

% =========================================================
% INDEX PAGE
% =========================================================
\newpage
\begin{center}
{\Huge \textbf{\textcolor{navyblue}{INDEX}}}\\[0.6cm]
\end{center}

\arrayrulecolor{navyblue}
\setlength{\arrayrulewidth}{1pt}
\renewcommand{\arraystretch}{1.4}

\begin{table}[h!]
\centering
\begin{tabular}{|c|p{7.5cm}|c|c|c|}
\hline
\rowcolor{navyblue}
\textcolor{white}{\textbf{Exp No.}} & \textcolor{white}{\textbf{Name of Experiment}} & \textcolor{white}{\textbf{Date}} & \textcolor{white}{\textbf{Page No.}} & \textcolor{white}{\textbf{Faculty Signature}} \\
\hline
1 & Area Of Triangle & & & \\ \hline
\rowcolor{navyblue!5}
2 & Largest Among Three Numbers & & & \\ \hline
3 & Palindrome Numbers & & & \\ \hline
\rowcolor{navyblue!5}
4 & Fibonacci Series & & & \\ \hline
5 & Vowels And Consonants & & & \\ \hline
\rowcolor{navyblue!5}
6 & Sorting & & & \\ \hline
7 & Maximum And Minimum Element Of A Tuple & & & \\ \hline
\rowcolor{navyblue!5}
8 & List Of Frequencies & & & \\ \hline
9 & Frequency Of Characters Of A File & & & \\ \hline
\rowcolor{navyblue!5}
10 & Exception Handling & & & \\ \hline
11 & Student Class & & & \\ \hline
\rowcolor{navyblue!5}
12 & Inheritance & & & \\ \hline
13 & Numpy & & & \\ \hline
\rowcolor{navyblue!5}
14 & Panda Module & & & \\ \hline
15 & Matplotlib & & & \\ \hline
\end{tabular}
\end{table}

\newpage
% =========================================================
% EXPERIMENT 1
% =========================================================
\noindent
\textbf{\large EXPERIMENT NO: 1} \hfill \textbf{\large Date: \underline{\hspace{3cm}}}

\vspace{0.2cm}
\begin{center}
{\Huge \textbf{\textcolor{navyblue}{AREA OF TRIANGLE}}}
\end{center}
\vspace{0.3cm}

\subsection*{\textit{AIM}}
Write a program to Find the area of a right angled triangle

\subsection*{\textit{OBJECTIVES}}
\begin{itemize}
  \item To understand fundamental user input operations and type casting in Python.
  \item To execute arithmetic multiplication and division calculations for geometry.
  \item To format and display computed numerical results to the console output.
\end{itemize}

\subsection*{\textit{ALGORITHM}}
\begin{description}
  \item[Step 1 :- Start] Initialize the execution environment to accept input values for base and height of the right angled triangle.
  \item[Step 2 :- Input Base Length] Prompt the user to enter the base length of the triangle and convert it to an integer.
  \item[Step 3 :- Input Triangle Height] Prompt the user to enter the height of the triangle and convert it to an integer.
  \item[Step 4 :- Compute Product] Multiply the base and height values together and store the result in variable c.
  \item[Step 5 :- Calculate Triangle Area] Divide the product c by two to calculate the area of the right angled triangle.
  \item[Step 6 :- Display Output Area] Output the calculated area value along with descriptive text to the standard console display window.
  \item[Step 7 :- Stop] Terminate the execution of the Python program successfully and release system memory resources.
\end{description}

\subsection*{\textit{PROGRAM}}
\begin{lstlisting}[style=pythonstyle]
b = int(input("Enter the base")) 
h = int(input("Enter the height")) 
c = b * h 
area = c / 2 
print("Area of right angle triangle is:", area, "sqcm")
\end{lstlisting}

\subsection*{\textit{SAMPLE OUTPUT}}
\begin{tcolorbox}[colback=black!85,colframe=navyblue,title=\textbf{Terminal Output},coltitle=white,fonttitle=\bfseries\small]
\small\ttfamily\color{green!30!white}
Enter the base: 5\\
Enter the height: 10\\
Area of right angle triangle is: 25.0 sqcm
\end{tcolorbox}

\subsection*{\textit{RESULT}}
The program to calculate the area of a right-angled triangle was successfully implemented and verified.

\newpage

% =========================================================
% EXPERIMENT 2
% =========================================================
\noindent
\textbf{\large EXPERIMENT NO: 2} \hfill \textbf{\large Date: \underline{\hspace{3cm}}}

\vspace{0.2cm}
\begin{center}
{\Huge \textbf{\textcolor{navyblue}{LARGEST AMONG THREE NUMBERS}}}
\end{center}
\vspace{0.3cm}

\subsection*{\textit{AIM}}
Write a program to Find the greatest among three numbers

\subsection*{\textit{OBJECTIVES}}
\begin{itemize}
  \item To implement multi-way branching using conditional control structures.
  \item To utilize relational and logical AND operators for comparison.
  \item To determine maximum numerical value among user inputs.
\end{itemize}

\subsection*{\textit{ALGORITHM}}
\begin{description}
  \item[Step 1 :- Start] Begin execution of the program to compare three user input numbers and find the largest value.
  \item[Step 2 :- Input First Number] Prompt the user to enter the first number and convert the input string to an integer.
  \item[Step 3 :- Input Second Number] Prompt the user to enter the second number and convert the input string to an integer.
  \item[Step 4 :- Input Third Number] Prompt the user to enter the third number and convert the input string to an integer.
  \item[Step 5 :- Evaluate First Number] Check if the first number is strictly greater than both second and third input numbers.
  \item[Step 6 :- Evaluate Second Number] Check if the second number is strictly greater than the third number if first check fails.
  \item[Step 7 :- Assign Third Number] Consider the third number as largest if both previous conditional checks evaluate to false.
  \item[Step 8 :- Display Largest Result] Print the identified largest number to the terminal display with an appropriate descriptive text message.
  \item[Step 9 :- Stop] Conclude the execution of the number comparison program and return control to the operating system.
\end{description}

\subsection*{\textit{PROGRAM}}
\begin{lstlisting}[style=pythonstyle]
a = int(input("Enter first number: ")) 
b = int(input("Enter second number: ")) 
c = int(input("Enter third number: ")) 
if a > b and a > c: 
    print("largest number is ", a) 
elif b > c: 
    print("the largest number is ", b) 
else: 
    print("the largest number is ", c)
\end{lstlisting}

\subsection*{\textit{SAMPLE OUTPUT}}
\begin{tcolorbox}[colback=black!85,colframe=navyblue,title=\textbf{Terminal Output},coltitle=white,fonttitle=\bfseries\small]
\small\ttfamily\color{green!30!white}
Enter first number: 10\\
Enter second number: 5\\
Enter third number: 2\\
largest number is 10
\end{tcolorbox}

\subsection*{\textit{RESULT}}
The program to find the largest among three numbers was successfully executed and verified.

\newpage

% =========================================================
% EXPERIMENT 3
% =========================================================
\noindent
\textbf{\large EXPERIMENT NO: 3} \hfill \textbf{\large Date: \underline{\hspace{3cm}}}

\vspace{0.2cm}
\begin{center}
{\Huge \textbf{\textcolor{navyblue}{PALINDROME NUMBERS}}}
\end{center}
\vspace{0.3cm}

\subsection*{\textit{AIM}}
Write a program to check whether a given number is palindrome or not

\subsection*{\textit{OBJECTIVES}}
\begin{itemize}
  \item To implement indefinite iteration using while loop structure.
  \item To perform digit extraction using modulus and integer floor division.
  \item To reverse integer values and verify numeric symmetry.
\end{itemize}

\subsection*{\textit{ALGORITHM}}
\begin{description}
  \item[Step 1 :- Start] Begin the program to check whether a given user input integer number is a palindrome or not.
  \item[Step 2 :- Read Input Number] Prompt the user to enter a number and store it in variable n and temporary variable temp.
  \item[Step 3 :- Initialize Reverse Counter] Set the initial value of the reversed number variable rev to zero before starting loop.
  \item[Step 4 :- Check Loop Condition] Execute a while loop as long as the temporary variable temp is strictly greater than zero.
  \item[Step 5 :- Extract Remainder Digit] Calculate the remainder digit by taking modulus of temporary variable temp by base value ten.
  \item[Step 6 :- Update Reversed Value] Multiply rev by ten and add extracted remainder digit to accumulate reversed numerical value.
  \item[Step 7 :- Truncate Temporary Digit] Perform integer floor division on variable temp by ten to remove the last processed digit.
  \item[Step 8 :- Compare Original Value] Compare the reversed number rev with original number n for exact numerical equality match.
  \item[Step 9 :- Print Palindrome Result] Print palindrome if numbers are equal, otherwise print not palindrome to console display window.
  \item[Step 10 :- Stop] Terminate the execution of the palindrome check Python script and clear runtime variables from memory.
\end{description}

\subsection*{\textit{PROGRAM}}
\begin{lstlisting}[style=pythonstyle]
n = int(input("enter a number")) 
rev = 0 
temp = n 
while temp > 0: 
    rem = temp % 10 
    rev = rev * 10 + rem 
    temp = temp // 10 

if rev == n: 
    print("palindrome") 
else: 
    print("not palindrome")
\end{lstlisting}

\subsection*{\textit{SAMPLE OUTPUT}}
\begin{tcolorbox}[colback=black!85,colframe=navyblue,title=\textbf{Terminal Output},coltitle=white,fonttitle=\bfseries\small]
\small\ttfamily\color{green!30!white}
enter a number: 121\\
palindrome
\end{tcolorbox}

\subsection*{\textit{RESULT}}
The palindrome number checking program was executed and verified for various input cases.

\newpage

% =========================================================
% EXPERIMENT 4
% =========================================================
\noindent
\textbf{\large EXPERIMENT NO: 4} \hfill \textbf{\large Date: \underline{\hspace{3cm}}}

\vspace{0.2cm}
\begin{center}
{\Huge \textbf{\textcolor{navyblue}{FIBONACCI SERIES}}}
\end{center}
\vspace{0.3cm}

\subsection*{\textit{AIM}}
Write a program to print Fibonacci series

\subsection*{\textit{OBJECTIVES}}
\begin{itemize}
  \item To understand sequence generation and iterative loop bounds.
  \item To manage variable swapping and accumulator updates during iterations.
  \item To display sequence terms up to user specified limit.
\end{itemize}

\subsection*{\textit{ALGORITHM}}
\begin{description}
  \item[Step 1 :- Start] Initialize the execution environment to compute and display a specified number of Fibonacci series terms.
  \item[Step 2 :- Read Terms Count] Prompt user to enter total number of terms and store integer value in variable n.
  \item[Step 3 :- Initialize Seed Values] Set first term variable a to zero and second term variable b to value one.
  \item[Step 4 :- Display Initial Seeds] Print the first two seed numbers a and b on separate lines to console.
  \item[Step 5 :- Execute Iterative Loop] Loop through index range starting from two up to specified total term count n.
  \item[Step 6 :- Calculate Next Term] Add current values of variables a and b to compute the next term c.
  \item[Step 7 :- Print Computed Term] Print the newly calculated Fibonacci term c to standard terminal console display screen.
  \item[Step 8 :- Shift Sequence Variables] Update variable a with value of b, and update variable b with new value c.
  \item[Step 9 :- Stop] Terminate the execution of Fibonacci series generation program and free allocated system memory resources.
\end{description}

\subsection*{\textit{PROGRAM}}
\begin{lstlisting}[style=pythonstyle]
n = int(input("enter your number of terms")) 
a = 0 
b = 1 
print(a) 
print(b) 
for i in range(2, n): 
    c = a + b 
    print(c)  
    a = b 
    b = c
\end{lstlisting}

\subsection*{\textit{SAMPLE OUTPUT}}
\begin{tcolorbox}[colback=black!85,colframe=navyblue,title=\textbf{Terminal Output},coltitle=white,fonttitle=\bfseries\small]
\small\ttfamily\color{green!30!white}
enter your number of terms: 7\\
0\\
1\\
1\\
2\\
3\\
5\\
8
\end{tcolorbox}

\subsection*{\textit{RESULT}}
The program to generate the Fibonacci series was implemented and output verified successfully.

\newpage

% =========================================================
% EXPERIMENT 5
% =========================================================
\noindent
\textbf{\large EXPERIMENT NO: 5} \hfill \textbf{\large Date: \underline{\hspace{3cm}}}

\vspace{0.2cm}
\begin{center}
{\Huge \textbf{\textcolor{navyblue}{VOWELS AND CONSONANTS}}}
\end{center}
\vspace{0.3cm}

\subsection*{\textit{AIM}}
Write a program to count number of vowels and consonants in a string

\subsection*{\textit{OBJECTIVES}}
\begin{itemize}
  \item To iterate over character elements of a text string.
  \item To perform string inspection using logical OR conditions and isalpha().
  \item To count distinct character classes in textual input.
\end{itemize}

\subsection*{\textit{ALGORITHM}}
\begin{description}
  \item[Step 1 :- Start] Begin the execution of string character parsing program to count total vowels and consonants.
  \item[Step 2 :- Input Text String] Prompt user to enter a string and store string value in variable s.
  \item[Step 3 :- Calculate String Length] Find total character count of string s using len function and store in l.
  \item[Step 4 :- Initialize Class Counters] Set vowel count variable v to zero and consonant count variable c to zero.
  \item[Step 5 :- Loop String Characters] Loop through character index range from index zero up to string length limit l.
  \item[Step 6 :- Check Vowel Character] Check if character at index i matches any lowercase or uppercase vowel letter.
  \item[Step 7 :- Increment Vowel Count] Add one to vowel counter v if character matches vowel comparison condition check.
  \item[Step 8 :- Check Consonant Character] Check if character is alphabetic letter using isalpha and increment consonant counter c.
  \item[Step 9 :- Display Summary Counts] Print total count of vowels and consonants to standard console terminal output window.
  \item[Step 10 :- Stop] Conclude execution of string parsing script and release string buffer memory back to system.
\end{description}

\subsection*{\textit{PROGRAM}}
\begin{lstlisting}[style=pythonstyle]
s = str(input("enter a string")) 
l = len(s) 
c = 0 
v = 0 
for i in range(l): 
    if s[i] == 'a' or s[i] == 'e' or s[i] == 'i' or s[i] == 'o' or s[i] == 'u' or \
       s[i] == 'A' or s[i] == 'E' or s[i] == 'I' or s[i] == 'O' or s[i] == 'U': 
        v = v + 1 
    elif s[i].isalpha(): 
        c = c + 1 
print("Number of vowels is ", v) 
print("Number of consonants is ", c)
\end{lstlisting}

\subsection*{\textit{SAMPLE OUTPUT}}
\begin{tcolorbox}[colback=black!85,colframe=navyblue,title=\textbf{Terminal Output},coltitle=white,fonttitle=\bfseries\small]
\small\ttfamily\color{green!30!white}
enter a string: Hello\\
Number of vowels is  2\\
Number of consonants is  3
\end{tcolorbox}

\subsection*{\textit{RESULT}}
The vowel and consonant counting program was successfully executed and output verified.

\newpage

% =========================================================
% EXPERIMENT 6
% =========================================================
\noindent
\textbf{\large EXPERIMENT NO: 6} \hfill \textbf{\large Date: \underline{\hspace{3cm}}}

\vspace{0.2cm}
\begin{center}
{\Huge \textbf{\textcolor{navyblue}{SORTING}}}
\end{center}
\vspace{0.3cm}

\subsection*{\textit{AIM}}
Write a program to Sort the elements of a list using function

\subsection*{\textit{OBJECTIVES}}
\begin{itemize}
  \item To define modular functions accepting list parameters in Python.
  \item To implement bubble sort algorithm using nested loops and element swapping.
  \item To return and print ordered data structures.
\end{itemize}

\subsection*{\textit{ALGORITHM}}
\begin{description}
  \item[Step 1 :- Start] Initialize the bubble sort program to rearrange list elements in non decreasing numerical order.
  \item[Step 2 :- Define Sort Function] Declare bubble sort function taking list argument arr to perform element sorting operations.
  \item[Step 3 :- Get List Length] Calculate total element count of input list arr and store integer length in n.
  \item[Step 4 :- Outer Pass Loop] Execute outer loop for variable i ranging from index zero to length n minus one.
  \item[Step 5 :- Inner Comparison Loop] Execute inner loop for variable j ranging from index zero to n minus i minus one.
  \item[Step 6 :- Compare Adjacent Values] Check if element at index j is strictly greater than adjacent element at j plus one.
  \item[Step 7 :- Swap Out Order Elements] Swap adjacent elements in place if left element is larger than right element.
  \item[Step 8 :- Return Sorted Array] Return ordered list arr from bubble sort function after completing all comparison passes.
  \item[Step 9 :- Call Function and Print] Call sorting function with sample list and print sorted result to terminal window.
  \item[Step 10 :- Stop] Conclude execution of sorting algorithm program and free list memory allocation from system.
\end{description}

\subsection*{\textit{PROGRAM}}
\begin{lstlisting}[style=pythonstyle]
def bubble_sort(arr): 
    n = len(arr) 
    for i in range(n): 
        for j in range(0, n - i - 1): 
            if arr[j] > arr[j + 1]: 
                arr[j], arr[j + 1] = arr[j + 1], arr[j] 
    return arr 
 
sample_list = [64, 34, 25, 12, 22, 11, 90] 
sorted_list = bubble_sort(sample_list) 
print(f"Sorted list in non-decreasing order: {sorted_list}")
\end{lstlisting}

\subsection*{\textit{SAMPLE OUTPUT}}
\begin{tcolorbox}[colback=black!85,colframe=navyblue,title=\textbf{Terminal Output},coltitle=white,fonttitle=\bfseries\small]
\small\ttfamily\color{green!30!white}
Sorted list in non-decreasing order: [11, 12, 22, 25, 34, 64, 90]
\end{tcolorbox}

\subsection*{\textit{RESULT}}
The list sorting function was successfully implemented using bubble sort and output verified.

\newpage

% =========================================================
% EXPERIMENT 7
% =========================================================
\noindent
\textbf{\large EXPERIMENT NO: 7} \hfill \textbf{\large Date: \underline{\hspace{3cm}}}

\vspace{0.2cm}
\begin{center}
{\Huge \textbf{\textcolor{navyblue}{MAXIMUM AND MINIMUM ELEMENT OF A TUPLE}}}
\end{center}
\vspace{0.3cm}

\subsection*{\textit{AIM}}
Write a program to find Maximum and minimum elements of a tuple using function

\subsection*{\textit{OBJECTIVES}}
\begin{itemize}
  \item To explore immutable sequence types (tuples) in Python.
  \item To traverse tuple elements and perform conditional minimum/maximum evaluation.
  \item To return tuple pairs containing multiple values from functions.
\end{itemize}

\subsection*{\textit{ALGORITHM}}
\begin{description}
  \item[Step 1 :- Start] Initialize program execution to find maximum and minimum elements of a tuple using function.
  \item[Step 2 :- Define Function] Declare function find max min accepting tuple parameter tup for processing element extremes.
  \item[Step 3 :- Handle Empty Input] Return None for both maximum and minimum if input tuple tup is empty.
  \item[Step 4 :- Set Initial Extremes] Set initial variables max element and min element to first element of tup.
  \item[Step 5 :- Loop Tuple Elements] Iterate through each numeric element num present inside input tuple tup in order.
  \item[Step 6 :- Update Maximum Value] Compare num with max element and update max element if num is strictly larger.
  \item[Step 7 :- Update Minimum Value] Compare num with min element and update min element if num is strictly smaller.
  \item[Step 8 :- Return Extreme Pair] Return tuple containing calculated max element and min element from custom function.
  \item[Step 9 :- Call and Display] Pass sample tuple to function and print maximum and minimum values to terminal screen.
  \item[Step 10 :- Stop] Conclude tuple analysis program execution successfully and clear tuple structures from memory.
\end{description}

\subsection*{\textit{PROGRAM}}
\begin{lstlisting}[style=pythonstyle]
def find_max_min(tup): 
    if not tup: 
        return None, None 
    max_element = tup[0] 
    min_element = tup[0] 
   
    for num in tup: 
        if num > max_element: 
            max_element = num 
        if num < min_element: 
            min_element = num 
    return max_element, min_element 
 
sample_tuple = (3, 1, 4, 1, 5, 9, 2, 6, 5, 3, 5) 
max_element, min_element = find_max_min(sample_tuple) 
print(f"The maximum element in the tuple is: {max_element}") 
print(f"The minimum element in the tuple is: {min_element}")
\end{lstlisting}

\subsection*{\textit{SAMPLE OUTPUT}}
\begin{tcolorbox}[colback=black!85,colframe=navyblue,title=\textbf{Terminal Output},coltitle=white,fonttitle=\bfseries\small]
\small\ttfamily\color{green!30!white}
The maximum element in the tuple is: 9\\
The minimum element in the tuple is: 1
\end{tcolorbox}

\subsection*{\textit{RESULT}}
The function to find maximum and minimum elements of a tuple was executed and verified.

\newpage

% =========================================================
% EXPERIMENT 8
% =========================================================
\noindent
\textbf{\large EXPERIMENT NO: 8} \hfill \textbf{\large Date: \underline{\hspace{3cm}}}

\vspace{0.2cm}
\begin{center}
{\Huge \textbf{\textcolor{navyblue}{LIST OF FREQUENCIES}}}
\end{center}
\vspace{0.3cm}

\subsection*{\textit{AIM}}
Write a program to create list of frequencies in non-increasing order using function

\subsection*{\textit{OBJECTIVES}}
\begin{itemize}
  \item To build frequency maps using Python dictionaries.
  \item To extract dictionary values into lists and implement custom sorting logic.
  \item To arrange frequency counts in descending numerical order.
\end{itemize}

\subsection*{\textit{ALGORITHM}}
\begin{description}
  \item[Step 1 :- Start] Initialize frequency distribution analyzer program to count list occurrences and sort in descending order.
  \item[Step 2 :- Define Frequency Function] Declare function frequency list accepting array parameter arr for calculating item occurrence counts.
  \item[Step 3 :- Check Empty List] Return empty list immediately if input list arr contains no numerical elements.
  \item[Step 4 :- Build Occurrence Counts] Iterate items in list arr, updating element counts inside dictionary frequency dict.
  \item[Step 5 :- Extract Frequency List] Convert dictionary count values into linear list frequencies for sorting processing.
  \item[Step 6 :- Outer Sorting Loop] Run outer loop over frequency list indices to begin descending bubble sort algorithm.
  \item[Step 7 :- Inner Sorting Loop] Run inner loop comparing current frequency against subsequent elements in frequencies list.
  \item[Step 8 :- Swap Descending Order] Swap elements if preceding frequency count is smaller than following frequency count.
  \item[Step 9 :- Return and Print] Return sorted frequency list from function and print result to standard output screen.
  \item[Step 10 :- Stop] Conclude execution of frequency counting script and release dictionary memory back to system.
\end{description}

\subsection*{\textit{PROGRAM}}
\begin{lstlisting}[style=pythonstyle]
def frequency_list(arr): 
    if not arr:  
        return [] 
    frequency_dict = {} 
 
    for item in arr: 
        if item in frequency_dict: 
            frequency_dict[item] += 1 
        else: 
            frequency_dict[item] = 1 
    frequencies = list(frequency_dict.values()) 
 
    for i in range(len(frequencies)): 
        for j in range(i + 1, len(frequencies)): 
            if frequencies[i] < frequencies[j]: 
                frequencies[i], frequencies[j] = frequencies[j], frequencies[i] 
 
    return frequencies 

sample_list = [3, 1, 4, 1, 5, 9, 2, 6, 5, 3, 5] 
frequencies = frequency_list(sample_list) 
print(f"List of frequencies in non-increasing order: {frequencies}")
\end{lstlisting}

\subsection*{\textit{SAMPLE OUTPUT}}
\begin{tcolorbox}[colback=black!85,colframe=navyblue,title=\textbf{Terminal Output},coltitle=white,fonttitle=\bfseries\small]
\small\ttfamily\color{green!30!white}
List of frequencies in non-increasing order: [3, 2, 2, 1, 1, 1, 1]
\end{tcolorbox}

\subsection*{\textit{RESULT}}
The frequency list computation program in non-increasing order was tested and verified.

\newpage

% =========================================================
% EXPERIMENT 9
% =========================================================
\noindent
\textbf{\large EXPERIMENT NO: 9} \hfill \textbf{\large Date: \underline{\hspace{3cm}}}

\vspace{0.2cm}
\begin{center}
{\Huge \textbf{\textcolor{navyblue}{FREQUENCY OF CHARACTERS OF A FILE}}}
\end{center}
\vspace{0.3cm}

\subsection*{\textit{AIM}}
Write a program to find Frequency of characters in a given file

\subsection*{\textit{OBJECTIVES}}
\begin{itemize}
  \item To manage file I/O operations using Python context managers (with open).
  \item To handle file stream exceptions like FileNotFoundError.
  \item To count character occurrences across file text content.
\end{itemize}

\subsection*{\textit{ALGORITHM}}
\begin{description}
  \item[Step 1 :- Start] Initialize character frequency counting program to process and count character occurrences in a file.
  \item[Step 2 :- Define File Function] Create function character frequency taking filename argument to read file and count characters.
  \item[Step 3 :- Open File Reader] Open file in read mode safely using open context manager and store content in text.
  \item[Step 4 :- Build Character Map] Iterate through each character char in text, updating counts in dictionary frequency dict.
  \item[Step 5 :- Exception Handling Guard] Intercept FileNotFoundError exception using try except block if target file does not exist.
  \item[Step 6 :- Create Sample File] Open file.txt in write mode, write string hello world, and close file handle.
  \item[Step 7 :- Call and Display] Invoke character frequency function with filename and print calculated frequency dictionary to console.
  \item[Step 8 :- Stop] Finish file operation execution context and release file stream resources back to operating system.
\end{description}

\subsection*{\textit{PROGRAM}}
\begin{lstlisting}[style=pythonstyle]
def character_frequency(filename): 
    try: 
        with open(filename, 'r') as file: 
            text = file.read() 
        frequency_dict = {} 
        for char in text: 
            if char in frequency_dict: 
                frequency_dict[char] += 1 
            else: 
                frequency_dict[char] = 1 
        return frequency_dict 
    except FileNotFoundError: 
        print(f"The file {filename} does not exist.") 
        return {} 

filename = 'file.txt' 
f = open(filename, 'w') 
f.write("hello world") 
f.close() 

frequencies = character_frequency(filename) 
print(f"Frequencies of characters in the file: {frequencies}")
\end{lstlisting}

\subsection*{\textit{SAMPLE OUTPUT}}
\begin{tcolorbox}[colback=black!85,colframe=navyblue,title=\textbf{Terminal Output},coltitle=white,fonttitle=\bfseries\small]
\small\ttfamily\color{green!30!white}
Frequencies of characters in the file: {'h': 1, 'e': 1, 'l': 3, 'o': 2, ' ': 1, 'w': 1, 'r': 1, 'd': 1}
\end{tcolorbox}

\subsection*{\textit{RESULT}}
The file character frequency analyzer was implemented and output verified successfully.

\newpage

% =========================================================
% EXPERIMENT 10
% =========================================================
\noindent
\textbf{\large EXPERIMENT NO: 10} \hfill \textbf{\large Date: \underline{\hspace{3cm}}}

\vspace{0.2cm}
\begin{center}
{\Huge \textbf{\textcolor{navyblue}{EXCEPTION HANDLING}}}
\end{center}
\vspace{0.3cm}

\subsection*{\textit{AIM}}
Write a program to handle exception in the case of voters age validity If age < 0 raise an exception

\subsection*{\textit{OBJECTIVES}}
\begin{itemize}
  \item To raise custom exceptions (ValueError) based on business logic.
  \item To wrap error prone statements inside try-except blocks.
  \item To intercept specific exception classes (ValueError, TypeError).
\end{itemize}

\subsection*{\textit{ALGORITHM}}
\begin{description}
  \item[Step 1 :- Start] Begin voter eligibility checking program with custom exception handling for negative age values.
  \item[Step 2 :- Define Validation Function] Declare function check voter age accepting numerical integer age parameter for validation evaluation.
  \item[Step 3 :- Check Negative Age] Raise ValueError exception with error message if input age is strictly less than zero.
  \item[Step 4 :- Evaluate Age Eligibility] Print not eligible if age is under eighteen, or eligible if age is eighteen or above.
  \item[Step 5 :- Read Input inside Try] Prompt user for integer age input inside try block and call check voter age function.
  \item[Step 6 :- Catch ValueError Exception] Catch raised ValueError exception and print descriptive error message to terminal console.
  \item[Step 7 :- Catch TypeError Exception] Catch TypeError exception if user provides invalid non integer input for age variable.
  \item[Step 8 :- Stop] Conclude exception handling program safely without crashing application runtime environment on invalid user inputs.
\end{description}

\subsection*{\textit{PROGRAM}}
\begin{lstlisting}[style=pythonstyle]
def check_voter_age(age): 
    if age < 0: 
        raise ValueError("Age cannot be negative.") 
    elif age < 18: 
        print("You are not eligible to vote.") 
    else: 
        print("You are eligible to vote.") 

try: 
    age = int(input("Enter your age: ")) 
    check_voter_age(age) 
except ValueError as e: 
    print(f"Error: {e}") 
except TypeError: 
    print("Please enter a valid integer for age.")
\end{lstlisting}

\subsection*{\textit{SAMPLE OUTPUT}}
\begin{tcolorbox}[colback=black!85,colframe=navyblue,title=\textbf{Terminal Output},coltitle=white,fonttitle=\bfseries\small]
\small\ttfamily\color{green!30!white}
Enter your age: -5\\
Error: Age cannot be negative.
\end{tcolorbox}

\subsection*{\textit{RESULT}}
The voter age exception handling program was executed and exception handling verified.

\newpage

% =========================================================
% EXPERIMENT 11
% =========================================================
\noindent
\textbf{\large EXPERIMENT NO: 11} \hfill \textbf{\large Date: \underline{\hspace{3cm}}}

\vspace{0.2cm}
\begin{center}
{\Huge \textbf{\textcolor{navyblue}{STUDENT CLASS}}}
\end{center}
\vspace{0.3cm}

\subsection*{\textit{AIM}}
Write a program to create a student class and create following methods: setMarks() - to update marks, getMarks() - to get individual marks, avgMarks() - to find average of marks, showResults() - to show results of each student

\subsection*{\textit{OBJECTIVES}}
\begin{itemize}
  \item To understand Object-Oriented Programming (OOP) concepts in Python.
  \item To initialize instance attributes using constructors (__init__).
  \item To define instance methods for data manipulation, query, and formatted summary display.
\end{itemize}

\subsection*{\textit{ALGORITHM}}
\begin{description}
  \item[Step 1 :- Start] Initialize student management application using object oriented programming concepts to model student scores.
  \item[Step 2 :- Define Student Class] Define Student class with constructor initializing name attribute and empty marks dictionary.
  \item[Step 3 :- Define setMarks Method] Define setMarks method taking subject and mark parameters to update student marks dictionary.
  \item[Step 4 :- Define getMarks Method] Define getMarks method to return individual subject score or default missing message string.
  \item[Step 5 :- Define avgMarks Method] Define avgMarks method calculating average by dividing sum of marks by total subjects.
  \item[Step 6 :- Define showResults Method] Define showResults method printing subject scores breakdown and computed average mark summary.
  \item[Step 7 :- Instantiate and Execute] Create Student object student1, set marks for subjects, and call class methods to display results.
  \item[Step 8 :- Stop] Conclude execution of object oriented student class program and destroy object instance from memory.
\end{description}

\subsection*{\textit{PROGRAM}}
\begin{lstlisting}[style=pythonstyle]
class Student: 
    def __init__(self, name): 
        self.name = name 
        self.marks = {} 
 
    def setMarks(self, subject, mark): 
        self.marks[subject] = mark 
 
    def getMarks(self, subject): 
        return self.marks.get(subject, "No marks available for this subject.") 
 
    def avgMarks(self): 
        if not self.marks: 
            return "No marks available." 
        total_marks = sum(self.marks.values()) 
        average_marks = total_marks / len(self.marks) 
        return average_marks 
 
    def showResults(self): 
        if not self.marks: 
            print(f"No marks available for {self.name}.") 
            return 
        print(f"Results for {self.name}:") 
        for subject, mark in self.marks.items(): 
            print(f"{subject}: {mark}") 
        average_marks = self.avgMarks() 
        print(f"Average Marks: {average_marks}") 

student1 = Student("John Doe") 
student1.setMarks("Math", 85) 
student1.setMarks("Science", 90) 
student1.setMarks("English", 88) 

print(student1.getMarks("Math")) 
print(student1.getMarks("History")) 
print(student1.avgMarks()) 
student1.showResults()
\end{lstlisting}

\subsection*{\textit{SAMPLE OUTPUT}}
\begin{tcolorbox}[colback=black!85,colframe=navyblue,title=\textbf{Terminal Output},coltitle=white,fonttitle=\bfseries\small]
\small\ttfamily\color{green!30!white}
85\\
No marks available for this subject.\\
87.66666666666667\\
Results for John Doe:\\
Math: 85\\
Science: 90\\
English: 88\\
Average Marks: 87.66666666666667
\end{tcolorbox}

\subsection*{\textit{RESULT}}
The Student class implementation with method encapsulation was verified successfully.

\newpage

% =========================================================
% EXPERIMENT 12
% =========================================================
\noindent
\textbf{\large EXPERIMENT NO: 12} \hfill \textbf{\large Date: \underline{\hspace{3cm}}}

\vspace{0.2cm}
\begin{center}
{\Huge \textbf{\textcolor{navyblue}{INHERITANCE}}}
\end{center}
\vspace{0.3cm}

\subsection*{\textit{AIM}}
Write a program to create a vehicle class with variables name, speed, mileage and methods setCapacity() and getFare(), create a class bus that inherits from the Vehicle class.

\subsection*{\textit{OBJECTIVES}}
\begin{itemize}
  \item To implement OOP inheritance between superclasses and subclasses.
  \item To leverage superclass methods using super().
  \item To demonstrate method overriding and specialized behavior in derived classes.
\end{itemize}

\subsection*{\textit{ALGORITHM}}
\begin{description}
  \item[Step 1 :- Start] Initialize inheritance program to model base Vehicle class and derived Bus child class behaviors.
  \item[Step 2 :- Define Vehicle Class] Define base Vehicle class with constructor setting name, speed, mileage, and initial capacity.
  \item[Step 3 :- Define Vehicle Methods] Add setCapacity method to store passenger limit and getFare method calculating base fare.
  \item[Step 4 :- Define Bus Subclass] Create child Bus class inheriting from Vehicle, overriding getFare method with additional charge.
  \item[Step 5 :- Instantiate Vehicle Object] Create Vehicle object vehicle1, set capacity to five, and display vehicle properties and fare.
  \item[Step 6 :- Instantiate Bus Object] Create Bus object bus1, set capacity to fifty, and display inherited properties and calculated fare.
  \item[Step 7 :- Stop] Conclude object oriented inheritance program execution and release class instances back to memory heap.
\end{description}

\subsection*{\textit{PROGRAM}}
\begin{lstlisting}[style=pythonstyle]
class Vehicle: 
    def __init__(self, name, speed, mileage): 
        self.name = name 
        self.speed = speed 
        self.mileage = mileage 
        self.capacity = 0 
 
    def setCapacity(self, capacity): 
        self.capacity = capacity 

    def getFare(self): 
        return self.capacity * 100 
 
class Bus(Vehicle): 
    def getFare(self): 
        base_fare = super().getFare() 
        return base_fare * 1.1 
 
vehicle1 = Vehicle("Car", 150, 15) 
vehicle1.setCapacity(5) 
print(f"Vehicle Name: {vehicle1.name}") 
print(f"Speed: {vehicle1.speed} km/h") 
print(f"Mileage: {vehicle1.mileage} km/l") 
print(f"Capacity: {vehicle1.capacity}") 
print(f"Fare: {vehicle1.getFare()}") 
 
bus1 = Bus("School Bus", 80, 10) 
bus1.setCapacity(50) 
print(f"\nBus Name: {bus1.name}") 
print(f"Speed: {bus1.speed} km/h") 
print(f"Mileage: {bus1.mileage} km/l") 
print(f"Capacity: {bus1.capacity}") 
print(f"Fare: {bus1.getFare()}")
\end{lstlisting}

\subsection*{\textit{SAMPLE OUTPUT}}
\begin{tcolorbox}[colback=black!85,colframe=navyblue,title=\textbf{Terminal Output},coltitle=white,fonttitle=\bfseries\small]
\small\ttfamily\color{green!30!white}
Vehicle Name: Car\\
Speed: 150 km/h\\
Mileage: 15 km/l\\
Capacity: 5\\
Fare: 500\\
\\
Bus Name: School Bus\\
Speed: 80 km/h\\
Mileage: 10 km/l\\
Capacity: 50\\
Fare: 5500.0
\end{tcolorbox}

\subsection*{\textit{RESULT}}
The OOP inheritance hierarchy between Vehicle and Bus classes was successfully demonstrated.

\newpage

% =========================================================
% EXPERIMENT 13
% =========================================================
\noindent
\textbf{\large EXPERIMENT NO: 13} \hfill \textbf{\large Date: \underline{\hspace{3cm}}}

\vspace{0.2cm}
\begin{center}
{\Huge \textbf{\textcolor{navyblue}{NUMPY}}}
\end{center}
\vspace{0.3cm}

\subsection*{\textit{AIM}}
Write a program to create a matrix using numpy and do the following operations: Perform matrix addition, subtraction and multiplication, Find transpose of matrix, Calculates the inverse of matrix, Calculates the determinant of matrix, Transforms a matrix into 1d array (flatten)

\subsection*{\textit{OBJECTIVES}}
\begin{itemize}
  \item To utilize NumPy library for numerical array processing.
  \item To perform fundamental matrix operations (addition, subtraction, dot product).
  \item To compute linear algebra functions (transpose, inverse, determinant, flattening).
\end{itemize}

\subsection*{\textit{ALGORITHM}}
\begin{description}
  \item[Step 1 :- Start] Initialize matrix computation module using NumPy library for matrix arithmetic and linear algebra operations.
  \item[Step 2 :- Import NumPy Package] Import numpy module under standard alias np to access array creation and matrix functions.
  \item[Step 3 :- Create Input Matrices] Create two 2D arrays matrix1 and matrix2 with 2x2 integer values for calculations.
  \item[Step 4 :- Compute Matrix Arithmetic] Perform element wise addition, subtraction, and dot product multiplication using NumPy built in functions.
  \item[Step 5 :- Calculate Transpose Matrix] Compute matrix transpose using np.transpose function to swap rows and columns of matrix1.
  \item[Step 6 :- Calculate Inverse and Determinant] Compute matrix inverse using np.linalg.inv and matrix determinant using np.linalg.det on matrix1.
  \item[Step 7 :- Flatten Matrix Array] Transform 2D array matrix1 into 1D array flattened using matrix flatten method call.
  \item[Step 8 :- Display Matrix Results] Print original matrices and calculated addition, subtraction, multiplication, transpose, inverse, determinant, and flattened outputs.
  \item[Step 9 :- Stop] Conclude NumPy matrix operations program execution and free array memory allocations from environment heap.
\end{description}

\subsection*{\textit{PROGRAM}}
\begin{lstlisting}[style=pythonstyle]
import numpy as np 
 
# Create matrices 
matrix1 = np.array([[1, 2], [3, 4]]) 
matrix2 = np.array([[5, 6], [7, 8]]) 
 
# Perform matrix addition 
addition = np.add(matrix1, matrix2) 
 
# Perform matrix subtraction 
subtraction = np.subtract(matrix1, matrix2) 
 
# Perform matrix multiplication 
multiplication = np.dot(matrix1, matrix2) 
 
# Find transpose of matrix1 
transpose = np.transpose(matrix1) 
 
# Calculate the inverse of matrix1 
inverse = np.linalg.inv(matrix1) 
 
# Calculate the determinant of matrix1 
determinant = np.linalg.det(matrix1) 
 
# Transform matrix1 into a 1D array (flatten) 
flattened = matrix1.flatten() 
 
# Display results 
print("Matrix 1:\n", matrix1) 
print("\nMatrix 2:\n", matrix2) 
print("\nMatrix Addition:\n", addition) 
print("\nMatrix Subtraction:\n", subtraction) 
print("\nMatrix Multiplication:\n", multiplication) 
print("\nTranspose of Matrix 1:\n", transpose) 
print("\nInverse of Matrix 1:\n", inverse) 
print("\nDeterminant of Matrix 1:\n", determinant) 
print("\nFlattened Matrix 1:\n", flattened)
\end{lstlisting}

\subsection*{\textit{SAMPLE OUTPUT}}
\begin{tcolorbox}[colback=black!85,colframe=navyblue,title=\textbf{Terminal Output},coltitle=white,fonttitle=\bfseries\small]
\small\ttfamily\color{green!30!white}
Matrix 1:\\
[[1 2]\\
 [3 4]]\\
\\
Matrix 2:\\
[[5 6]\\
 [7 8]]\\
\\
Matrix Addition:\\
[[ 6  8]\\
 [10 12]]\\
\\
Matrix Subtraction:\\
[[-4 -4]\\
 [-4 -4]]\\
\\
Matrix Multiplication:\\
[[19 22]\\
 [43 50]]\\
\\
Transpose of Matrix 1:\\
[[1 3]\\
 [2 4]]\\
\\
Inverse of Matrix 1:\\
[[-2.   1. ]\\
 [ 1.5 -0.5]]\\
\\
Determinant of Matrix 1:\\
-2.0\\
\\
Flattened Matrix 1:\\
[1 2 3 4]
\end{tcolorbox}

\subsection*{\textit{RESULT}}
The matrix operations program using NumPy library was executed and results verified.

\newpage

% =========================================================
% EXPERIMENT 14
% =========================================================
\noindent
\textbf{\large EXPERIMENT NO: 14} \hfill \textbf{\large Date: \underline{\hspace{3cm}}}

\vspace{0.2cm}
\begin{center}
{\Huge \textbf{\textcolor{navyblue}{PANDA MODULE}}}
\end{center}
\vspace{0.3cm}

\subsection*{\textit{AIM}}
Write a python program to read a .csv file using panda module and print the first and last five lines of a file

\subsection*{\textit{OBJECTIVES}}
\begin{itemize}
  \item To leverage Pandas library for structured data manipulation.
  \item To load CSV datasets into DataFrames using read_csv.
  \item To inspect head and tail slices of dataset files.
\end{itemize}

\subsection*{\textit{ALGORITHM}}
\begin{description}
  \item[Step 1 :- Start] Initialize data processing program using Pandas module to read CSV file records and preview data lines.
  \item[Step 2 :- Import Pandas Module] Import pandas library under alias pd to utilize DataFrame loading and inspection utility functions.
  \item[Step 3 :- Read CSV File] Read external CSV dataset file using pd.read_csv function and store data in DataFrame df.
  \item[Step 4 :- Extract Head Rows] Call head method on DataFrame df to retrieve the first five rows of CSV dataset.
  \item[Step 5 :- Extract Tail Rows] Call tail method on DataFrame df to retrieve the last five rows of CSV dataset.
  \item[Step 6 :- Display Preview Slices] Print first five lines and last five lines of dataset to standard terminal console display.
  \item[Step 7 :- Stop] Terminate execution of Pandas file processing script and clear DataFrame structures from system memory.
\end{description}

\subsection*{\textit{PROGRAM}}
\begin{lstlisting}[style=pythonstyle]
import pandas as pd 
 
filename = 'example.csv' 
df = pd.read_csv(filename) 
 
print("First five lines of the file:") 
print(df.head()) 
 
print("\nLast five lines of the file:") 
print(df.tail())
\end{lstlisting}

\subsection*{\textit{SAMPLE OUTPUT}}
\begin{tcolorbox}[colback=black!85,colframe=navyblue,title=\textbf{Terminal Output},coltitle=white,fonttitle=\bfseries\small]
\small\ttfamily\color{green!30!white}
First five lines of the file:\\
      Name  Age         City\\
0    Alice   30     New York\\
1      Bob   25  Los Angeles\\
2  Charlie   35      Chicago\\
3    David   40      Houston\\
4      Eve   28      Phoenix\\
\\
Last five lines of the file:\\
     Name  Age         City\\
5   Frank   33  Philadelphia\\
6   Grace   45   San Antonio\\
7  Hannah   27     San Diego\\
8    Ivan   50        Dallas\\
9    Judy   29      San Jose
\end{tcolorbox}

\subsection*{\textit{RESULT}}
The Pandas CSV data reading and head/tail preview program was verified successfully.

\newpage

% =========================================================
% EXPERIMENT 15
% =========================================================
\noindent
\textbf{\large EXPERIMENT NO: 15} \hfill \textbf{\large Date: \underline{\hspace{3cm}}}

\vspace{0.2cm}
\begin{center}
{\Huge \textbf{\textcolor{navyblue}{MATPLOTLIB}}}
\end{center}
\vspace{0.3cm}

\subsection*{\textit{AIM}}
Write a program to draw geometric shapes using matplotlib

\subsection*{\textit{OBJECTIVES}}
\begin{itemize}
  \item To initialize Matplotlib figure and subplot axis objects.
  \item To render geometric patch primitives (Rectangle, Circle, Polygon).
  \item To format plot aesthetics including labels, title, grid, aspect ratio, and show display.
\end{itemize}

\subsection*{\textit{ALGORITHM}}
\begin{description}
  \item[Step 1 :- Start] Initialize graphic visualization program using Matplotlib module to render geometric shapes on plot canvas.
  \item[Step 2 :- Import Plotting Modules] Import matplotlib.pyplot module under alias plt and numpy module under alias np for drawing.
  \item[Step 3 :- Create Figure Canvas] Create figure and subplots axis objects using plt.subplots function to define plot drawing canvas.
  \item[Step 4 :- Add Rectangle Shape] Create Rectangle patch specifying position, width, height, color, and alpha transparency, adding to axis.
  \item[Step 5 :- Add Circle Shape] Create Circle patch specifying center position, radius, color, and transparency, adding to plot axis.
  \item[Step 6 :- Add Triangle Polygon] Create Polygon patch with NumPy array triangle coordinates, adding red triangle patch to axis.
  \item[Step 7 :- Configure Axis Limits] Set x and y axis limits from zero to one and set equal aspect ratio scaling.
  \item[Step 8 :- Display Plot Figure] Set grid, axis labels, title, and display rendered figure window using plt.show function.
  \item[Step 9 :- Stop] Close graphic visualization window and terminate program execution releasing plot resources back to system.
\end{description}

\subsection*{\textit{PROGRAM}}
\begin{lstlisting}[style=pythonstyle]
import matplotlib.pyplot as plt 
import numpy as np 
 
fig, ax = plt.subplots() 
 
rectangle = plt.Rectangle((0.2, 0.5), 0.4, 0.3, color='blue', alpha=0.3) 
ax.add_patch(rectangle) 
 
circle = plt.Circle((0.8, 0.2), 0.1, color='green', alpha=0.5) 
ax.add_patch(circle) 
 
triangle_points = np.array([[0.6, 0.8], [0.8, 0.8], [0.7, 0.9], [0.6, 0.8]]) 
triangle = plt.Polygon(triangle_points, closed=True, color='red', alpha=0.4) 
ax.add_patch(triangle) 
 
ax.set_xlim(0, 1) 
ax.set_ylim(0, 1) 
ax.set_aspect('equal', adjustable='box') 
 
ax.grid(True) 
ax.set_xlabel('X-axis') 
ax.set_ylabel('Y-axis') 
ax.set_title('Geometric Shapes') 
 
plt.show()
\end{lstlisting}

\subsection*{\textit{SAMPLE OUTPUT}}
\begin{tcolorbox}[colback=black!85,colframe=navyblue,title=\textbf{Terminal Output},coltitle=white,fonttitle=\bfseries\small]
\small\ttfamily\color{green!30!white}
[Plot Window Displayed: Rendered Blue Rectangle, Green Circle, and Red Triangle shapes on grid axes]
\end{tcolorbox}

\subsection*{\textit{RESULT}}
The Matplotlib geometric shape rendering program was successfully executed and visualized.

\newpage

\end{document}
